\section{The Pen Wars: Venice, England, Barclay (1606--1610)}
\label{sec:penwars}

\subsection{The Venetian interdict}

The first great quarrel of Paul~V's pontificate was
jurisdictional. The Republic of Venice had subjected clerics to
its civil courts in certain causes and restricted the Church's
acquisition of real property; the pope laid Venice under interdict
in 1606. The paper war that followed set the Servite Paolo Sarpi
and the theologians of the Republic against the theologians of the
Holy See: ``a war of pamphlets,'' the 1907 Catholic Encyclopedia
records, ``in which the part of the Republic was sustained by John
Marsiglio and an apostate monk named Paolo Sarpi, and that of the
Holy See by Bellarmine and Baronius'' (E/K, reported; the
encyclopedia's polemical epithet for Sarpi is its own). Bellarmine
contributed Italian tracts against the Venetian
theologians---``two or three little Italian books against the
doctors of Venice,'' as his autobiography inventories them
(D\"ollinger--Reusch, 43; A). The interdict was resolved by French
mediation in 1607 without retraction on either side; the episode
mattered for Bellarmine chiefly as the first live test, before all
Europe, of the ecclesiology of \work{De Summo Pontifice}---and as
the opening round of a decade in which his opponents were no
longer professors but states.

\subsection{The Oath of Allegiance and the duel with James~I}

In the aftermath of the Gunpowder Plot, the English statute of
1606 imposed on Catholics an Oath of Allegiance. Its sting was not
the profession of civil loyalty but a clause requiring the juror
to ``abhorre, detest and abiure, as impious and Hereticall, this
damnable doctrine and position, That Princes which be
excommunicated or depriued by the Pope, may be deposed or
murthered by their Subiects or any other whatsoeuer'' (oath text
as reprinted within James's \work{Triplici nodo, triplex cuneus}
in the 1616 \work{Workes}, McIlwain ed.; N). Paul~V forbade the
oath in two breves, dated (in the Roman style the king's book
reproduces) the tenth of the calends of October 1606---22
September---and in the calends of September 1607; and on 28
September 1607 Bellarmine wrote to the archpriest George
Blackwell, who had taken the oath, condemning his course and
declaring the oath contrary to the Catholic faith (McIlwain,
introduction and reprinted texts; N/K). The king himself
(anonymously at first) answered the two breves and the cardinal's
letter together in \work{Triplici nodo, triplex cuneus: or an
Apologie for the Oath of Allegiance} (1607). Bellarmine replied in 1608 under the
name of his chaplain Matteo Torti---\work{Responsio Matthaei
Torti}---arguing that the oath as worded made the juror deny, as
``impious and heretical,'' a doctrine that popes had acted on and
theologians defended for centuries, and that loyalty could have
been secured without it. James reissued his \work{Apologie} over
his own name in 1609 with the long \work{Premonition} to the
princes of Europe; Bellarmine answered with the \work{Apologia
pro responsione sua} (1609), now over his own name (CE 1907 for
the sequence of titles; N for the English texts in McIlwain's
edition).

The exchange escalated from pamphlet to European incident: the
\work{Premonition} of 1609 was a summons to all Christian princes
against the deposing power as such, and the quarrel fed directly
into the Gallican crisis in France, where the crown's lawyers
regarded the indirect power as scarcely less seditious than the
direct. The 1907 Catholic Encyclopedia draws the summary line: for
following the via media of the indirect power he was ``condemned
in 1590 as too much of a Regalist'' and later as too much of a
Papalist.

\subsection{Against Barclay: the indirect power stated}

William Barclay, a Scottish jurist in France, had died leaving a
treatise, \work{De potestate Papae}, denying the pope any power,
direct or indirect, in temporal matters. Bellarmine's reply,
\work{Tractatus de potestate Summi Pontificis in rebus
temporalibus, adversus Gulielmum Barclay} (1610), is the fullest
statement of his mature position: the pope has no ordinary civil
jurisdiction over princes; but in virtue of his spiritual charge
he may, \emph{indirectly} and for spiritual ends, dispose of
temporal things---including, in the extreme case, absolving
subjects and deposing rulers whose rule destroys souls. The
Parlement of Paris condemned the book (CE 1907, ``branded by the
Regalist Parlement of Paris''); his autobiography lists it among
the works of the period along with ``an apologetic book against
the King of England'' and a reply to the English theologian Roger
Widdrington (D\"ollinger--Reusch, 43; A).

\witnesslimit{Interpretive ceiling: Bellarmine's indirect-power
doctrine was a theory of the relation of two jurisdictions within
Christendom, argued from medieval precedent; it is not a modern
doctrine of religious liberty, nor of the separation of church
and state, and this study does not modernize it. Equally, the
doctrine's coercive implications---deposition, absolution of
subjects---are part of its content and are stated as such. Both
its seventeenth-century enemies (regalist and papalist) and its
later admirers have enlisted it for causes beyond its terms; the
reception section returns to this.}
